\section{Related work}
\label{sec:related}

\paragraph{Hallucination in tool use.}
Tool-calling benchmarks measure whether a model can call the right function with the right arguments
\citep{patil2024gorilla,qin2024toolllm,patil2025bfcl}, and the current leaderboard adds an explicit category in
which no available tool applies and the model is expected to abstain \citep{patil2025bfcl}. These evaluate
capability. The question here is different: given that a model sometimes fails, can the failure be detected from
the model's own state before the call executes. \citet{healy2026internal} first showed that it can, with a probe
on residual-stream states at the function-name, argument and delimiter positions, and that probe is the
residual-stream baseline we compare against and extend. Concurrently, \citet{yeats2026calls} train linear
probes on the hidden states of eighteen models on the same benchmark and find that the probes catch
argument errors of the right type and wrong value and generalise to unseen error kinds. That work establishes
breadth on the residual tier; it has no attention tier, no surface floor and no held-out-tool split, which are
the three things that decide the comparison this paper makes.

\paragraph{Probing the residual stream.}
Supervised probes on hidden states detect factual error \citep{azaria2023internal,marks2024geometry}, predict
hallucination risk from the query alone \citep{ji2024internal}, and stand in for sampling-based estimates at
the cost of a single pass \citep{kossen2024sep}. \citet{orgad2025know} showed that the signal concentrates at
the exact-answer tokens rather than being spread over the response, which is the observation the token-role
probe applies to a call's name, arguments and delimiter. The family has been developed in several further
directions: covariance and eigenvalue statistics of the states
\citep{chen2024inside}, token-level probes trained on entity-level annotations that operate in real time on
long-form generation \citep{obeso2025realtime}, cross-layer dynamics of the residual stream
\citep{icrprobe2025}, spectral analysis of hidden-layer trajectories \citep{li2025hsad}, and probes over factual
claims \citep{bhatnagar2026halt}. Process reward models verify reasoning steps with a similar supervised
signal \citep{lightman2024verify}. All of these require residual-stream access. Our contribution to this family
is not a new probe but the observation that the token-role probe remains the strongest single detector we
measured, and that the correct comparison for it is an attention-only method at matched protocol.

\paragraph{Reading attention.}
Attention is the tensor practitioners are most likely to be able to read, and it has been used to decide
cache eviction \citep{zhang2023h2o}, to detect hallucination through the ratio of attention paid to context
against generation \citep{chuang2024lookback}, to aggregate multi-layer attention into a single map
\citep{abnar2020quantifying}, and to score consistency of attention across depth
\citep{badash2026layers}. The spectral line treats a layer as a weighted graph or kernel:
\citet{sriramanan2024llmcheck} read eigenvalues of attention kernel maps and of hidden-state covariance in a
single pass, \citet{noel2025gsp} introduced Laplacian diagnostics of the attention graph, and
\citet{binkowski2025lapeigvals} obtained the strongest attention-only results to date from top-$k$ Laplacian
values; the same group has since read hallucination from attention sinks \citep{binkowski2026sinks}, a
per-head signal we do not evaluate here. \citet{dahlem2026transport} prove that every transpose-invariant
spectral diagnostic of attention is blind to the direction of attention flow and evaluate spectral detectors
under length control; every spectrum in this paper is of a symmetrised graph and is therefore
orientation-blind in their sense, which bounds what the attention tier can see and is why our length floor
sits in every table. Two features of this literature motivate our theory section. First, nearly all of it averages over
heads before computing anything, which Section~\ref{sec:theory} shows is the decisive step. Second, part of it
computes quantities that require residual states as a graph signal while presenting them as properties of
attention; we separate the two and tag every feature by the tensors it consumes.

\paragraph{Sampling and output-side signals.}
Semantic entropy and self-consistency detect hallucination by sampling several continuations
\citep{farquhar2024semantic,manakul2023selfcheckgpt}. They are strong but they multiply inference cost and they
require the freedom to sample, which an agent about to execute an action does not have. We therefore restrict
this study to a single forward pass and include the cheapest output-side signal, mean token log-probability, as a
baseline rather than as a competitor.

\paragraph{Position of this work.}
The two internals families were developed independently, and the open question was which to deploy. We answer it
as a frontier over access rather than a ranking, and in the course of answering it we find that the interesting
axis is not residual versus attention but averaged versus per head. The negative half of the result is as
informative as the positive half: at matched protocol our per-head detector does not beat
\citet{binkowski2025lapeigvals}, and Section~\ref{sec:theory:lapeig} says why the two should be close despite
computing different functionals.
